%%
%% This is file `sample-manuscript.tex',
%% generated with the docstrip utility.
%%
%% The original source files were:
%%
%% samples.dtx  (with options: `manuscript')
%% 
%% IMPORTANT NOTICE:
%% 
%% For the copyright see the source file.
%% 
%% Any modified versions of this file must be renamed
%% with new filenames distinct from sample-manuscript.tex.
%% 
%% For distribution of the original source see the terms
%% for copying and modification in the file samples.dtx.
%% 
%% This generated file may be distributed as long as the
%% original source files, as listed above, are part of the
%% same distribution. (The sources need not necessarily be
%% in the same archive or directory.)
%%
%% Commands for TeXCount
%TC:macro \cite [option:text,text]
%TC:macro \citep [option:text,text]
%TC:macro \citet [option:text,text]
%TC:envir table 0 1
%TC:envir table* 0 1
%TC:envir tabular [ignore] word
%TC:envir displaymath 0 word
%TC:envir math 0 word
%TC:envir comment 0 0
%%
%%
%% The first command in your LaTeX source must be the \documentclass command.
%%%% Small single column format, used for CIE, CSUR, DTRAP, JACM, JDIQ, JEA, JERIC, JETC, PACMCGIT, TAAS, TACCESS, TACO, TALG, TALLIP (formerly TALIP), TCPS, TDSCI, TEAC, TECS, TELO, THRI, TIIS, TIOT, TISSEC, TIST, TKDD, TMIS, TOCE, TOCHI, TOCL, TOCS, TOCT, TODAES, TODS, TOIS, TOIT, TOMACS, TOMM (formerly TOMCCAP), TOMPECS, TOMS, TOPC, TOPLAS, TOPS, TOS, TOSEM, TOSN, TQC, TRETS, TSAS, TSC, TSLP, TWEB.
% \documentclass[acmsmall]{acmart}

%%%% Large single column format, used for IMWUT, JOCCH, PACMPL, POMACS, TAP, PACMHCI
% \documentclass[acmlarge,screen]{acmart}

%%%% Large double column format, used for TOG
% \documentclass[acmtog, authorversion]{acmart}

%%%% Generic manuscript mode, required for submission
%%%% and peer review
\documentclass[manuscript,screen,review]{acmart}
%% Fonts used in the template cannot be substituted; margin 
%% adjustments are not allowed.
%%
%% \BibTeX command to typeset BibTeX logo in the docs
\AtBeginDocument{%
  \providecommand\BibTeX{{%
    \normalfont B\kern-0.5em{\scshape i\kern-0.25em b}\kern-0.8em\TeX}}}

%% Rights management information.  This information is sent to you
%% when you complete the rights form.  These commands have SAMPLE
%% values in them; it is your responsibility as an author to replace
%% the commands and values with those provided to you when you
%% complete the rights form.
\setcopyright{acmcopyright}
\copyrightyear{2018}
\acmYear{2018}
\acmDOI{XXXXXXX.XXXXXXX}

%% These commands are for a PROCEEDINGS abstract or paper.
\acmConference[Conference acronym 'XX]{Make sure to enter the correct
  conference title from your rights confirmation emai}{June 03--05,
  2018}{Woodstock, NY}
%
%  Uncomment \acmBooktitle if th title of the proceedings is different
%  from ``Proceedings of ...''!
%
\acmBooktitle{Woodstock '18: ACM Symposium on Neural Gaze Detection,
 June 03--05, 2018, Woodstock, NY} 
\acmPrice{15.00}
\acmISBN{978-1-4503-XXXX-X/18/06}


%%
%% Submission ID.
%% Use this when submitting an article to a sponsored event. You'll
%% receive a unique submission ID from the organizers
%% of the event, and this ID should be used as the parameter to this command.
%%\acmSubmissionID{123-A56-BU3}

%%
%% For managing citations, it is recommended to use bibliography
%% files in BibTeX format.
%%
%% You can then either use BibTeX with the ACM-Reference-Format style,
%% or BibLaTeX with the acmnumeric or acmauthoryear sytles, that include
%% support for advanced citation of software artefact from the
%% biblatex-software package, also separately available on CTAN.
%%
%% Look at the sample-*-biblatex.tex files for templates showcasing
%% the biblatex styles.
%%

%%
%% The majority of ACM publications use numbered citations and
%% references.  The command \citestyle{authoryear} switches to the
%% "author year" style.
%%
%% If you are preparing content for an event
%% sponsored by ACM SIGGRAPH, you must use the "author year" style of
%% citations and references.
%% Uncommenting
%% the next command will enable that style.
%%\citestyle{acmauthoryear}

%%
%% end of the preamble, start of the body of the document source.
\begin{document}

\title{Task 2  Numerical Methods}
\author{Carl-Emil Rohde \\ kvb125@alumni.ku.dk}
\date{September 2023}

\maketitle

\section*{Problem 1}

\textbf{No Solution}
\\
In this case, the equations will be inconsistent, meaning that they can't all be true at the same time.
\\
$$
\begin{aligned}
& x+y+z=1 \\
& 2 x+2 y+2 z=2 \\
& 3 x+3 y+3 z=5
\end{aligned}
$$
\\
\textbf{One Solution Only}
\\
In this case, the equations will be independent and intersect at exactly one point in $\mathbb{R}^3$.
\\
$$
\begin{aligned}
& x+y+z=1 \\
& x-y+z=3 \\
& 2 x+3 y-z=1
\end{aligned}
$$
\\
Solving these equations yields the unique solution $(x, y, z)=(1,-1,1)$.
\\
\\
\textbf{Infinitely Many Solutions}
\\
In this case, the equations are dependent and represent planes that intersect along a line or coincide.
\\
$$
\begin{aligned}
& x+y+z=1 \\
& 2 x+2 y+2 z=2 \\
& 3 x+3 y+3 z=3
\end{aligned}
$$
\\
All three equations are multiples of each other and therefore represent the same plane in $\mathbb{R}^3$ . Any point on that plane is a solution, so there are infinitely many solutions.
\section*{Problem 2}
\subsection*{4.2.30}
Crout's method allows me to compute $L$ and $U$ as follows:
\\
For $i=1,2,3, \ldots, n$ (where $n$ is the dimension of the square matrix $A$ ):
\\
\textbf{1. }$u_{i i}=1$
\\
\textbf{2.} $l_{i j}=a_{i j}-\sum_{k=1}^{i-1} l_{i k} u_{k j}$ (for $j=i, i+1, \ldots, n$ )
\\
\textbf{3.} $u_{i j}=\frac{1}{l_{i i}}\left(a_{i j}-\sum_{k=1}^{i-1} l_{i k} u_{k j}\right)$ (for $j=i+1, i+2, \ldots, n$ )
\\
\\
Let's start with the given matrix $A$ :
$$
A=\left[\begin{array}{rrr}
3 & 0 & 1 \\
0 & -1 & 3 \\
1 & 3 & 0
\end{array}\right]
$$
\text { I then initialize } L \text { and } U \text { to zero matrices of the same dimensions as } A \text {. } and calculate $L$ and $U$:
\\
$$
\begin{aligned}
& \textbf{i=1:} \\
& \cdot u_{11}=1 \\
& \cdot l_{11}=3-0=3 \\
& \cdot l_{21}=0-0=0 \\
& \cdot l_{31}=1-0=1 \\
& \cdot u_{12}=(0-0) / 3=0 \\
& \cdot u_{13}=(1-0) / 3=1 / 3 \\
& \textbf{i=2}: \\
& \cdot u_{22}=1 \\
& \cdot l_{22}=-1-0=-1 \\
& \cdot l_{32}=3-0=3 \\
& \cdot u_{23}=(3-0) /(-1)=-3 \\
& \textbf{i=3}: \\
& \cdot u_{33}=1 \\
& \cdot l_{33}=0-(1 \times-3)=3
\end{aligned}
$$
After this i get:
$$
\begin{aligned}
& L=\left[\begin{array}{rrr}
3 & 0 & 0 \\
0 & -1 & 0 \\
1 & 3 & 3
\end{array}\right] \\
& U=\left[\begin{array}{rrr}
1 & 0 & 1 / 3 \\
0 & 1 & -3 \\
0 & 0 & 1
\end{array}\right]
\end{aligned}
$$
\subsection*{4.2.31}
The given matrix $A$ is:
$$
A=\left[\begin{array}{ll}
1 & 2 \\
2 & 5
\end{array}\right]
$$
and the $L$ matrix is defined as follows:
$$
L=\left[\begin{array}{ll}
l_{11} & 0 \\
l_{21} & l_{22}
\end{array}\right]
$$
I will use the equations $A=L L^T$ to find $l_{11}, l_{21}$, and $l_{22}$:
$$
l_{11}^2=1 \Rightarrow l_{11}=1
$$
$$
l_{21} \times l_{11}=2 \Rightarrow l_{21}=2
$$
$$
l_{21}^2+l_{22}^2=5 \Rightarrow l_{22}=\sqrt{5-4} \Rightarrow l_{22}=\sqrt{1} \Rightarrow l_{22}=1
$$
The $L$ matrix then becomes:
$$
L=\left[\begin{array}{ll}
1 & 0 \\
2 & 1
\end{array}\right]
$$
Now $A=L L^T$ :
$$
A=\left[\begin{array}{ll}
1 & 0 \\
2 & 1
\end{array}\right]\left[\begin{array}{ll}
1 & 2 \\
0 & 1
\end{array}\right]=\left[\begin{array}{ll}
1 & 2 \\
2 & 5
\end{array}\right]
$$
So, the Cholesky decomposition of $A$ into $L L^T$ results in $L$ being the $2 \times 2$ matrix $\left[\begin{array}{ll}1 & 0 \\ 2 & 1\end{array}\right]$
\subsection*{4.2.32}
To find the $L L^T$-factorization of $A$, where $L$ is a lower triangular matrix with positive diagonal elements, I set up $L$ as:
$$
L=\left[\begin{array}{ccc}
l_{11} & 0 & 0 \\
l_{21} & l_{22} & 0 \\
l_{31} & l_{32} & l_{33}
\end{array}\right]
$$
We then have $A=L L^T$. I can now write the equations as follows:
$$
\begin{aligned}
& l_{11}^2=4 \Rightarrow l_{11}=2 \\
& l_{21} \times l_{11}=\frac{1}{2} \Rightarrow l_{21}=\frac{1}{4} \\
& l_{31} \times l_{11}=1 \Rightarrow l_{31}=\frac{1}{2} \\
& l_{21}^2+l_{22}^2=\frac{17}{16} \Rightarrow l_{22}=\sqrt{\frac{17}{16}-\frac{1}{16}} \Rightarrow l_{22}=\sqrt{\frac{16}{16}} \Rightarrow l_{22}=1 \\
& l_{31} \times l_{21}+l_{32} \times l_{22}=\frac{1}{4} \Rightarrow l_{32}=\frac{1}{4}-\frac{1}{8} \Rightarrow l_{32}=\frac{1}{8} \\
& l_{31}^2+l_{32}^2+l_{33}^2=\frac{33}{64} \Rightarrow l_{33}=\sqrt{\frac{33}{64}-\frac{4}{64}-\frac{1}{64}} \Rightarrow l_{33}=\sqrt{\frac{28}{64}} \Rightarrow l_{33}=\frac{1}{2}
\end{aligned}
$$
After solving these equations i get the matrix $L$:
$$
L=\left[\begin{array}{ccc}
2 & 0 & 0 \\
\frac{1}{4} & 1 & 0 \\
\frac{1}{2} & \frac{1}{8} & \frac{1}{2}
\end{array}\right]
$$
So, the $L L^T$-factorization of $A$ with $L$ having positive diagonal elements results in the above $L$ matrix.
\subsection*{4.4.16}
To find the condition number of a matrix $A$ using the infinity matrix norm, I can use the formula:
$$
\operatorname{cond}(A)=\|A\|_{\infty} \times\left\|A^{-1}\right\|_{\infty}
$$
Here, $A$ is the given matrix
$$
A=\left[\begin{array}{cc}
7 & 8 \\
9 & 10
\end{array}\right]
$$
First, I will find $\|A\|_{\infty}$ :
The infinity norm of $A$ is the maximum row sum:
$$
\|A\|_{\infty}=\max (7+8,9+10)=\max (15,19)=19
$$
Next, we find $A^{-1}$ using Cramer's Rule.
The determinant of $A$, denoted $\operatorname{det}(A)$, is:
$$
\operatorname{det}(A)=(7 \times 10)-(8 \times 9)=70-72=-2
$$
Using Cramer's Rule, the inverse $A^{-1}$ can be found as:
$$
A^{-1}=\frac{1}{\operatorname{det}(A)}\left[\begin{array}{cc}
10 & -8 \\
-9 & 7
\end{array}\right]=\left[\begin{array}{cc}
-5 & 4 \\
4.5 & -3.5
\end{array}\right]
$$
Now I will find $\left\|A^{-1}\right\|_{\infty}$ :
The infinity norm of $A^{-1}$ is the maximum row sum:
$$
\left\|A^{-1}\right\|_{\infty}=\max (-5+4,4.5+3.5)=\max (-1,8)=8
$$
Finally, the condition number is:
$$
\operatorname{cond}(A)=19 \times 8=152
$$
So the condition number of the matrix $A$ using the infinity norm is 152 .

\subsection*{Problem 4}
The Neumann Series provides a way to find the inverse of a matrix 
$A$ as:
$$
A^{-1}=I+(I-A)+(I-A)^2+(I-A)^3+\ldots
$$
The series will converge if $\|I-A\|<1$.
\\
To find $A$, I subtract $B$ from the identity matrix $I$ :
$$
A=I-B=\left(\begin{array}{lll}
1 & 0 & 0 \\
0 & 1 & 0 \\
0 & 0 & 1
\end{array}\right)-\left(\begin{array}{ccc}
0.5 & -0.2 & -0.2 \\
0.1 & 0.7 & -e \\
0.1 & -0.3 & 0.5
\end{array}\right)=\left(\begin{array}{ccc}
0.5 & 0.2 & 0.2 \\
-0.1 & 0.3 & e \\
-0.1 & 0.3 & 0.5
\end{array}\right)
$$
To determine for what values of $e$ the Neumann series converges, I will find the infinity matrix norm $\|A\|_{\infty}$ and determine when $\|A\|_{\infty}<1$.
\\
The infinity matrix norm of a matrix $A$ is defined as
$$
\|A\|_{\infty}=\max _i\left(\sum_j\left|a_{i j}\right|\right)
$$
For matrix $A$, the sum of the absolute values of the rows are:
$$
\begin{aligned}
& 0.5+0.2+0.2=0.9 \\
& 0.1+0.3+|e|=0.4+|e| \\
& 0.1+0.3+0.5=0.9
\end{aligned}
$$
The largest sum here is $\max (0.9,0.9,0.4+|e|)$.
For the Neumann series to converge, this value must be smaller than 1 . Therefore,
$$
\max (0.9,0.9,0.4+|e|)<1
$$
This inequality is satisfied for $|e|<0.6$.
So, for $e$ in the interval $-0.6<e<0.6$, the Neumann Series will converge to the inverse of $B$.
\subsection*{Problem 5}
The condition number $\kappa(A)$ of a matrix $A$ with respect to a given norm is defined as:
$$
\kappa(A)=\|A\| \times\left\|A^{-1}\right\|
$$
Firstly, I will find $A^{-1}$. Since $A$ is a diagonal matrix, $A^{-1}$ is simply the diagonal matrix where each diagonal entry $a_{i i}$ is replaced by $1 / a_{i i}$.
$$
A^{-1}=\left(\begin{array}{ccc}
1 & 0 & 0 \\
0 & \frac{1}{1} & 0 \\
0 & 0 & 1
\end{array}\right) \times\left(\begin{array}{ccc}
1 & 0 & 0 \\
0 & \frac{1}{1} & -\frac{e}{1} \\
0 & 0 & 1
\end{array}\right)=\left(\begin{array}{ccc}
1 & 0 & 0 \\
0 & 1 & -e \\
0 & 0 & 1
\end{array}\right)
$$
Now, we find the infinity norms $\|A\|_{\infty}$ and $\left\|A^{-1}\right\|_{\infty}$.
The infinity norm of a matrix is given by the maximum absolute row sum. 
\\
For $A$, each row sum is $|1|,|1+e|,|1|$. Therefore, $\|A\|_{\infty}=\max (1,1+|e|, 1)=1+|e|$.
\\
\\
Similarly, for $A^{-1}$, each row sum is $|1|,|1+| e \|,|1|$. Therefore, $\left\|A^{-1}\right\|_{\infty}=\max (1,1+$ $|e|, 1)=1+|e|$.
So, the condition number $\kappa(A)$ as a function of $e$ is:
$$
\kappa(A)=(1+|e|) \times(1+|e|)=(1+|e|)^2
$$


\end{document}
\endinput
%%
%% End of file `sample-authordraft.tex'.
